\documentclass[%
    font=nexus,
    aspectratio=169,
    language=english
]{tubsCD/tubscd}

\usepackage{hyperref}

\begin{document}

\begin{frame}[plain]
\titlepage
\end{frame}

\begin{frame}{Outline}
  \tableofcontents
\end{frame}

\part{Introduction}

\begin{frame}[plain]
  \partpage
\end{frame}

\section{Background}

\begin{frame}{What is XSS?}
  In \textbf{Cross-Site Scripting (XSS)}, an attacker tries to inject JavaScript into a webpage \citep{lekies2017scriptgadgets}.
  \vspace{0.5em}
  \begin{exampleblock}{Classic payload}
    \texttt{<script>alert(1)</script>}
  \end{exampleblock}
  \vspace{0.3em}
  XSS has ranked among the most common web vulnerabilities for years.
\end{frame}

\begin{frame}{Motivation}
  Defenses block \textbf{obvious executable code} \citep{lekies2017scriptgadgets}:
  \texttt{<script>}, event handlers, \texttt{javascript:} URLs.
  \vspace{0.6em}
  \textbf{Second line:} sanitizers, WAFs, browser filters, CSP.
  \vspace{0.6em}
  Modern apps also load \textbf{trusted JavaScript} from frameworks and libraries.
\end{frame}

\begin{frame}[plain]
  \centering
  \vfill
  \vspace{1em}
  {\Huge Can an attacker reuse this trusted code\\
  instead of injecting new JavaScript?\par}
  \vfill
  {\Large\bfseries Answer:\par}
  \vspace{0.8em}
  {\LARGE Yes, this is the idea behind \textbf{Script Gadgets}\par}
  \vfill
\end{frame}

\begin{frame}{Four Classes of XSS Mitigations}
  Per \cite{lekies2017scriptgadgets} (Section~2.3):
  \begin{enumerate}
    \item \textbf{HTML sanitizers} (DOMPurify, Google Closure)
    \item \textbf{Browser XSS filters} (Chrome, Edge, NoScript)
    \item \textbf{Web Application Firewalls} (ModSecurity CRS)
    \item \textbf{Content Security Policy (CSP)}
  \end{enumerate}
  \vspace{0.4em}
  Three underlying strategies:
  request filtering, response sanitization, code filtering.
  \vspace{0.3em}
  \small Most of these look mainly for \textbf{directly executable} patterns -- not for abuse of existing site code.
\end{frame}

\begin{frame}{But What If No Script Is Injected?}
  \begin{block}{Key question}
    What if the attacker does \textbf{not} inject a \texttt{<script>} tag at all?
  \end{block}
  \vspace{0.3em}
  \begin{enumerate}
    \item Classic XSS blocked by mitigations
    \item Attacker injects \textbf{harmless-looking HTML}
    \item Existing website JavaScript processes it
    \item A \textbf{Script Gadget} causes code execution
  \end{enumerate}
  \vspace{0.4em}
  \small
  Most XSS defenses assume: stop injected executable JavaScript.
  Script gadgets show this assumption is \textbf{incomplete} when the site's own code does the dangerous work.
\end{frame}

\part{Script Gadgets}

\begin{frame}[plain]
  \partpage
\end{frame}

\section{Concept}

\begin{frame}{Definition: Script Gadget}
  \begin{block}{Script Gadget \citep{lekies2017scriptgadgets} (Section~3.3)}
    \textbf{Legitimate} JavaScript already present on the website, which can be triggered by attacker-controlled but \textbf{seemingly harmless HTML}.
  \end{block}
  \begin{itemize}
    \item Not injected by the attacker (already in app or library code).
    \item Processes attacker data in an \textbf{unsafe} way once the gadget path activates.
    \item Introduces \textbf{code-reuse attacks} on the web (analogy: return-to-libc).
  \end{itemize}
\end{frame}

\begin{frame}{Benign HTML}
  Mitigations block directly executable code but often allow \citep{lekies2017scriptgadgets} (Section~3.1):
  \begin{itemize}
    \item No \texttt{<script>}, no inline handlers
    \item No \texttt{javascript:}/\texttt{data:} URLs in \texttt{src}/\texttt{href}
    \item No executable tags such as \texttt{<link rel=import>}, \texttt{<meta>}, \texttt{<style>}
  \end{itemize}
  \begin{exampleblock}{Listing 1 -- benign for filters}
    \small\texttt{<div class="greeting"><b>Hello</b> world!</div>}
  \end{exampleblock}
\end{frame}

\begin{frame}{DOM Selectors}
  Benign markup alone does not execute code (Section~3.2):
  \begin{itemize}
    \item Legitimate code reads the DOM (\texttt{getElementById}, \texttt{querySelector}, jQuery).
    \item Paper example: elements with \textbf{tooltip} attributes are decorated.
    \item Injected markup \textbf{matches} selectors $\rightarrow$ gadget path activates.
  \end{itemize}
\end{frame}

\begin{frame}{Attack Outline}
  Four steps \citep{lekies2017scriptgadgets} (Section~3.4):
  \begin{enumerate}
    \item \textbf{Injection} -- benign HTML (stored/reflected XSS)
    \item \textbf{Mitigation} -- no direct script detected, markup remains
    \item \textbf{Gadget} -- trusted code transforms the payload
    \item \textbf{Execution} -- attacker-controlled JavaScript runs
  \end{enumerate}
\end{frame}

\begin{frame}{Gadget Types}
  Classification (Section~3.5):
  \begin{itemize}
    \item \textbf{String manipulation} -- e.g.\ \texttt{data-*} $\rightarrow$ framework attributes (AngularJS, Polymer)
    \item \textbf{Element construction} -- new \texttt{<script>} (jQuery \texttt{\$.globalEval})
    \item \textbf{Function creation} -- \texttt{new Function(...)}; often combined with \texttt{unsafe-eval} CSP
  \end{itemize}
  Complex attacks \textbf{chain} multiple gadgets.
\end{frame}

\part{Bypassing Mitigations}

\begin{frame}[plain]
  \partpage
\end{frame}

\section{Paper Section 4}

\begin{frame}{Bypass Strategy (Overview)}
  Gadgets use \textbf{different HTML tags/attributes} than classic XSS payloads \citep{lekies2017scriptgadgets}:
  \begin{itemize}
    \item Request filtering: only ``evil'' patterns, not gadget paths
    \item Response sanitization: whitelist keeps markup; library makes it dangerous
    \item Code filtering (CSP): execution via \textbf{trusted} code
  \end{itemize}
\end{frame}

\begin{frame}{Bypass: Request Filtering}
  WAFs, NoScript (Section~4.2):
  \begin{itemize}
    \item Detect \texttt{<script>}, \texttt{onerror}, \texttt{eval}, \texttt{document.cookie}, \ldots
    \item Gadgets use \textbf{allowed} attributes (\texttt{data-html}, \texttt{data-bind}, \ldots)
    \item Frameworks split exploits (Knockout example, Listing~13): harmless parts together = attack
  \end{itemize}
\end{frame}

\begin{frame}{Bypass: Response Sanitization}
  HTML sanitizers, IE/Edge XSS filter (Section~4.3):
  \begin{itemize}
    \item Sanitizer: only whitelisted attributes (\texttt{class}, \texttt{id}) -- gadget abuses these
    \item jQuery Mobile + DOMPurify (Listing~14): \texttt{data-role=popup} injects comment $\rightarrow$ escape $\rightarrow$ XSS
    \item Gadgets are ``safe'' by definition until library logic activates them
  \end{itemize}
\end{frame}

\begin{frame}{Bypass: Code Filtering (CSP)}
  CSP and XSS Auditor (Section~4.4):
  \begin{itemize}
    \item \texttt{unsafe-eval}: gadgets using \texttt{eval}/templates (e.g.\ Underscore)
    \item \texttt{strict-dynamic}: trusted libs create new \texttt{<script>} elements (jQuery)
    \item Nonce/whitelist CSP: expression parsers in AngularJS, Aurelia, Polymer
    \item \textbf{13/16} frameworks studied: bypass of \texttt{strict-dynamic}
  \end{itemize}
\end{frame}

\begin{frame}{Empirical Results}
  \begin{itemize}
    \item 16 JS libraries/frameworks analyzed manually (jQuery, AngularJS, Vue, Bootstrap, \ldots)
    \item Gadgets in \textbf{almost all}; React had no usable gadgets
    \item Automated study: Alexa Top~5000, $\sim$650\,000 URLs
    \item \textbf{19.88\%} of domains: verified gadgets (conservative lower bound)
    \item 4\,352\,491 sink executions with DOM-sourced data measured
  \end{itemize}
  \small \cite{lekies2017scriptgadgets}
\end{frame}

\part{Case Study}

\begin{frame}[plain]
  \partpage
\end{frame}

\section{Course Material}

\begin{frame}{Example: Bootstrap Tooltip}
  Supplementary course example \cite{friendsDemo2026} -- illustrates the paper:
  \begin{itemize}
    \item \textbf{Stored XSS:} unsanitized HTML in profile/tooltip \texttt{title}
    \item \textbf{Gadget:} Bootstrap renders tooltip HTML (trusted library)
    \item \textbf{Bypass:} \texttt{sanitize: false} disables internal allowlist protection
    \item Trigger on hover; exfiltration e.g.\ via \texttt{<img onerror>} + \texttt{fetch}
  \end{itemize}
  \vspace{0.3em}
  \small Maps to Paper Sections~3--4; use the course lab app for a live demo if required.
\end{frame}

\part{Conclusion}

\begin{frame}[plain]
  \partpage
\end{frame}

\section{Wrap-up}

\begin{frame}{Countermeasures}
  \begin{itemize}
    \item Fix the flaw: context-aware encoding/sanitization \textbf{before} the DOM
    \item Do not disable library defaults (\texttt{sanitize: true})
    \item \textbf{CSP} helps but is not enough alone: review \texttt{strict-dynamic} and \texttt{unsafe-eval}
    \item Defense in depth: sanitizer + CSP + safe API usage
    \item PoC collection: \url{https://github.com/google/security-research-pocs}
  \end{itemize}
\end{frame}

\begin{frame}{Takeaways}
  \begin{alertblock}{Key message \citep{lekies2017scriptgadgets}}
    The assumption ``injected code $\neq$ benign markup'' is false in modern web apps.
  \end{alertblock}
  \begin{itemize}
    \item Script gadgets bypass all four mitigation classes.
    \item Trusted frameworks can be abused for XSS execution.
    \item Secure development is still required -- mitigations do not replace it.
  \end{itemize}
\end{frame}

\begin{frame}{Questions?}
  \centering
  \Large Thank you for your attention!

  \vspace{1em}
  \normalsize
  \end{frame}

\begin{frame}[allowframebreaks]{References}
    \bibliographystyle{apalike}
    \bibliography{bibliography.bib}
\end{frame}

\end{document}
